% Pollis Domain Packages: StochasticVolatility.jl
% Compile from this folder with: julia -e 'using PollisTeX; display(compile("StochasticVolatility.tex"))'
% The values come from values.jl, whose output is in ../benchmark/vignette_values.txt, and from the stored files of ../benchmark named in the text.
% The look of a Pollis method vignette (the fonts and colours of the Pollis New Methods documents).
% Set \vignetteauthor, \vignettedate and, if not CC BY-SA 4.0, \vignettelicense after % The look of a Pollis method vignette (the fonts and colours of the Pollis New Methods documents).
% Set \vignetteauthor, \vignettedate and, if not CC BY-SA 4.0, \vignettelicense after % The look of a Pollis method vignette (the fonts and colours of the Pollis New Methods documents).
% Set \vignetteauthor, \vignettedate and, if not CC BY-SA 4.0, \vignettelicense after \input{vignette_style}.
\documentclass[10pt,a4paper]{article}

\usepackage[a4paper,left=18mm,right=18mm,top=15mm,bottom=17mm,footskip=8mm]{geometry}
\usepackage{iftex}
\ifPDFTeX\usepackage[T1]{fontenc}\fi
\usepackage[default,semibold]{sourcesanspro}
\usepackage{textcomp}
\usepackage[italic]{mathastext}
\usepackage{amsmath}
\usepackage{microtype}
\usepackage[table]{xcolor}
\usepackage{tikz}
\usepackage[most]{tcolorbox}
\usepackage{fontawesome5}
\usepackage{tabularx,booktabs}
\usepackage{enumitem}
\usepackage{fancyhdr}
\usepackage[hidelinks]{hyperref}

% ---------------------------------------------------------------- colours (as in the brochure)
\definecolor{navy}{HTML}{1C3A5E}
\definecolor{pblue}{HTML}{3B78B5}
\definecolor{amber}{HTML}{F2B33D}
\definecolor{deepamber}{HTML}{D9902A}
\definecolor{pale}{HTML}{EEF3F9}
\definecolor{ink}{HTML}{2A2F36}
\definecolor{logonavy}{HTML}{173D63}
\definecolor{logodark}{HTML}{C98A1E}
\definecolor{logoamber}{HTML}{F2B446}
\definecolor{logocream}{HTML}{FEEFD2}
\definecolor{soft}{HTML}{5E6773}
\definecolor{agree}{HTML}{2F7D4F}
\definecolor{differ}{HTML}{B23B3B}
\color{ink}

\setlength{\parindent}{0pt}
\setlength{\parskip}{4pt}
\setlist{nosep,leftmargin=*,itemsep=2pt}
\renewcommand{\arraystretch}{1.15}

\pagestyle{fancy}
\fancyhf{}
\renewcommand{\headrulewidth}{0pt}
\fancyfoot[L]{\footnotesize\color{soft}\textbf{Pollis}\enspace New Methods\enspace·\enspace\nmtheme}
\fancyfoot[R]{\footnotesize\color{soft}\thepage}
\newcommand{\nmtheme}{}
\newcommand{\vignetteauthor}{}
\newcommand{\vignettedate}{\today}
\newcommand{\vignettelicense}{\href{https://creativecommons.org/licenses/by-sa/4.0/}{CC BY-SA 4.0}}

% The Pollis logo, redrawn so it prints sharply at any size
\newcommand{\pollislogo}[1]{%
  \begin{tikzpicture}[x=#1/512,y=-#1/512,baseline=(current bounding box.center)]
    \fill[logonavy] (0,0) rectangle (512,512);
    \fill[logodark] (36,422) -- (104,256) -- (172,422) -- cycle;
    \fill[logoamber] (104,422) -- (203,180) -- (302,422) -- cycle;
    \fill[logocream] (203,422) -- (339,89) -- (476,422) -- cycle;
  \end{tikzpicture}}

% Top of the first page: logo, name; then the kicker, title and a one-line promise
\newcommand{\nmhead}[3]{%
  \renewcommand{\nmtheme}{#1}%
  \thispagestyle{fancy}%
  \begin{minipage}[c]{13mm}\pollislogo{11mm}\end{minipage}%
  \begin{minipage}[c]{\dimexpr\linewidth-13mm\relax}
    {\color{navy}\Large\bfseries Pollis}\enspace{\color{soft}New Methods}\hfill
    {\color{deepamber}\small\bfseries\MakeUppercase{#1}}
  \end{minipage}\par\vspace{7mm}
  {\color{navy}\fontsize{22}{26}\selectfont\bfseries #2}\par\vspace{4pt}
  {\color{soft}\large #3}\par\vspace{5mm}}

\newcommand{\nmsection}[1]{\par\vspace{5pt}{\color{navy}\large\bfseries #1}\par\vspace{1pt}}
\newcommand{\nmsub}[1]{\par\vspace{2pt}{\color{deepamber}\bfseries #1}\par}
\newcommand{\gt}{\,{\color{deepamber}\bfseries ›}\,}
\newcommand{\menupath}[1]{{\bfseries\color{navy}#1}}
\newcommand{\fn}[1]{\texttt{\small #1}}
% Citations: \nmcite{key}{Breiman (2001)} links to the entry \nmref{key}{...} in the References
\newcommand{\nmcite}[2]{\hyperlink{ref:#1}{\color{pblue}#2}}
\newcommand{\nmref}[2]{\raisebox{1.1em}[0pt][0pt]{\hypertarget{ref:#1}{}}#2\par}
\newcommand{\nmdoi}[1]{\href{https://doi.org/#1}{\color{pblue}doi:#1}}
% Where the package is: its name and a link (its GitHub repository, once published)
\newcommand{\package}[2]{\href{#2}{\fn{#1.jl}}}

% Key facts at the top: source, missing in Julia, package, result
\newcommand{\factsheet}[4]{%
  \begin{tcolorbox}[enhanced, colback=pale, colframe=pale, boxrule=0pt, sharp corners,
    left=7pt, right=7pt, top=5pt, bottom=5pt, fontupper=\small]
  \begin{tabularx}{\linewidth}{@{}>{\leavevmode\bfseries\color{navy}}p{27mm}>{\leavevmode}X@{}}
  Tested against & #1\\
  Missing in Julia & #2\\
  Package & #3\\
  Result & #4
  \end{tabularx}
  \end{tcolorbox}}

% Boxes
\newtcolorbox{pbox}[2][]{enhanced, colback=pale, colframe=amber, boxrule=0pt,
  leftrule=3pt, sharp corners, left=7pt, right=7pt, top=6pt, bottom=6pt,
  fontupper=\small\raggedright, before upper={{\color{navy}\normalsize\bfseries #2}\par\smallskip}, #1}
\newtcolorbox{nbox}[2][]{enhanced, colback=navy, colframe=navy, boxrule=0pt,
  sharp corners, left=8pt, right=8pt, top=7pt, bottom=7pt, coltext=white,
  fontupper=\small\raggedright, before upper={{\color{amber}\normalsize\bfseries #2}\par\smallskip}, #1}

% Value tables: published beside Julia
\newenvironment{valuetable}[1]{%
  \par\small\tabularx{\linewidth}{#1}\toprule}{\bottomrule\endtabularx\par}
\newcommand{\ok}{{\color{agree}\faCheck}}
\newcommand{\no}{{\color{differ}\faTimes}}

% End of the document: credit and licence
\newcommand{\nmcredit}{%
  \par\vfill
  \begin{minipage}[b]{0.78\linewidth}\footnotesize\color{soft}
  A Pollis method, tested against the published tables on the published data.\par
  \textcopyright\ \vignetteauthor\enspace·\enspace Licensed under \vignettelicense\enspace·\enspace\vignettedate
  \end{minipage}\hfill
  \begin{minipage}[b]{0.18\linewidth}\raggedleft\pollislogo{9mm}\end{minipage}}
.
\documentclass[10pt,a4paper]{article}

\usepackage[a4paper,left=18mm,right=18mm,top=15mm,bottom=17mm,footskip=8mm]{geometry}
\usepackage{iftex}
\ifPDFTeX\usepackage[T1]{fontenc}\fi
\usepackage[default,semibold]{sourcesanspro}
\usepackage{textcomp}
\usepackage[italic]{mathastext}
\usepackage{amsmath}
\usepackage{microtype}
\usepackage[table]{xcolor}
\usepackage{tikz}
\usepackage[most]{tcolorbox}
\usepackage{fontawesome5}
\usepackage{tabularx,booktabs}
\usepackage{enumitem}
\usepackage{fancyhdr}
\usepackage[hidelinks]{hyperref}

% ---------------------------------------------------------------- colours (as in the brochure)
\definecolor{navy}{HTML}{1C3A5E}
\definecolor{pblue}{HTML}{3B78B5}
\definecolor{amber}{HTML}{F2B33D}
\definecolor{deepamber}{HTML}{D9902A}
\definecolor{pale}{HTML}{EEF3F9}
\definecolor{ink}{HTML}{2A2F36}
\definecolor{logonavy}{HTML}{173D63}
\definecolor{logodark}{HTML}{C98A1E}
\definecolor{logoamber}{HTML}{F2B446}
\definecolor{logocream}{HTML}{FEEFD2}
\definecolor{soft}{HTML}{5E6773}
\definecolor{agree}{HTML}{2F7D4F}
\definecolor{differ}{HTML}{B23B3B}
\color{ink}

\setlength{\parindent}{0pt}
\setlength{\parskip}{4pt}
\setlist{nosep,leftmargin=*,itemsep=2pt}
\renewcommand{\arraystretch}{1.15}

\pagestyle{fancy}
\fancyhf{}
\renewcommand{\headrulewidth}{0pt}
\fancyfoot[L]{\footnotesize\color{soft}\textbf{Pollis}\enspace New Methods\enspace·\enspace\nmtheme}
\fancyfoot[R]{\footnotesize\color{soft}\thepage}
\newcommand{\nmtheme}{}
\newcommand{\vignetteauthor}{}
\newcommand{\vignettedate}{\today}
\newcommand{\vignettelicense}{\href{https://creativecommons.org/licenses/by-sa/4.0/}{CC BY-SA 4.0}}

% The Pollis logo, redrawn so it prints sharply at any size
\newcommand{\pollislogo}[1]{%
  \begin{tikzpicture}[x=#1/512,y=-#1/512,baseline=(current bounding box.center)]
    \fill[logonavy] (0,0) rectangle (512,512);
    \fill[logodark] (36,422) -- (104,256) -- (172,422) -- cycle;
    \fill[logoamber] (104,422) -- (203,180) -- (302,422) -- cycle;
    \fill[logocream] (203,422) -- (339,89) -- (476,422) -- cycle;
  \end{tikzpicture}}

% Top of the first page: logo, name; then the kicker, title and a one-line promise
\newcommand{\nmhead}[3]{%
  \renewcommand{\nmtheme}{#1}%
  \thispagestyle{fancy}%
  \begin{minipage}[c]{13mm}\pollislogo{11mm}\end{minipage}%
  \begin{minipage}[c]{\dimexpr\linewidth-13mm\relax}
    {\color{navy}\Large\bfseries Pollis}\enspace{\color{soft}New Methods}\hfill
    {\color{deepamber}\small\bfseries\MakeUppercase{#1}}
  \end{minipage}\par\vspace{7mm}
  {\color{navy}\fontsize{22}{26}\selectfont\bfseries #2}\par\vspace{4pt}
  {\color{soft}\large #3}\par\vspace{5mm}}

\newcommand{\nmsection}[1]{\par\vspace{5pt}{\color{navy}\large\bfseries #1}\par\vspace{1pt}}
\newcommand{\nmsub}[1]{\par\vspace{2pt}{\color{deepamber}\bfseries #1}\par}
\newcommand{\gt}{\,{\color{deepamber}\bfseries ›}\,}
\newcommand{\menupath}[1]{{\bfseries\color{navy}#1}}
\newcommand{\fn}[1]{\texttt{\small #1}}
% Citations: \nmcite{key}{Breiman (2001)} links to the entry \nmref{key}{...} in the References
\newcommand{\nmcite}[2]{\hyperlink{ref:#1}{\color{pblue}#2}}
\newcommand{\nmref}[2]{\raisebox{1.1em}[0pt][0pt]{\hypertarget{ref:#1}{}}#2\par}
\newcommand{\nmdoi}[1]{\href{https://doi.org/#1}{\color{pblue}doi:#1}}
% Where the package is: its name and a link (its GitHub repository, once published)
\newcommand{\package}[2]{\href{#2}{\fn{#1.jl}}}

% Key facts at the top: source, missing in Julia, package, result
\newcommand{\factsheet}[4]{%
  \begin{tcolorbox}[enhanced, colback=pale, colframe=pale, boxrule=0pt, sharp corners,
    left=7pt, right=7pt, top=5pt, bottom=5pt, fontupper=\small]
  \begin{tabularx}{\linewidth}{@{}>{\leavevmode\bfseries\color{navy}}p{27mm}>{\leavevmode}X@{}}
  Tested against & #1\\
  Missing in Julia & #2\\
  Package & #3\\
  Result & #4
  \end{tabularx}
  \end{tcolorbox}}

% Boxes
\newtcolorbox{pbox}[2][]{enhanced, colback=pale, colframe=amber, boxrule=0pt,
  leftrule=3pt, sharp corners, left=7pt, right=7pt, top=6pt, bottom=6pt,
  fontupper=\small\raggedright, before upper={{\color{navy}\normalsize\bfseries #2}\par\smallskip}, #1}
\newtcolorbox{nbox}[2][]{enhanced, colback=navy, colframe=navy, boxrule=0pt,
  sharp corners, left=8pt, right=8pt, top=7pt, bottom=7pt, coltext=white,
  fontupper=\small\raggedright, before upper={{\color{amber}\normalsize\bfseries #2}\par\smallskip}, #1}

% Value tables: published beside Julia
\newenvironment{valuetable}[1]{%
  \par\small\tabularx{\linewidth}{#1}\toprule}{\bottomrule\endtabularx\par}
\newcommand{\ok}{{\color{agree}\faCheck}}
\newcommand{\no}{{\color{differ}\faTimes}}

% End of the document: credit and licence
\newcommand{\nmcredit}{%
  \par\vfill
  \begin{minipage}[b]{0.78\linewidth}\footnotesize\color{soft}
  A Pollis method, tested against the published tables on the published data.\par
  \textcopyright\ \vignetteauthor\enspace·\enspace Licensed under \vignettelicense\enspace·\enspace\vignettedate
  \end{minipage}\hfill
  \begin{minipage}[b]{0.18\linewidth}\raggedleft\pollislogo{9mm}\end{minipage}}
.
\documentclass[10pt,a4paper]{article}

\usepackage[a4paper,left=18mm,right=18mm,top=15mm,bottom=17mm,footskip=8mm]{geometry}
\usepackage{iftex}
\ifPDFTeX\usepackage[T1]{fontenc}\fi
\usepackage[default,semibold]{sourcesanspro}
\usepackage{textcomp}
\usepackage[italic]{mathastext}
\usepackage{amsmath}
\usepackage{microtype}
\usepackage[table]{xcolor}
\usepackage{tikz}
\usepackage[most]{tcolorbox}
\usepackage{fontawesome5}
\usepackage{tabularx,booktabs}
\usepackage{enumitem}
\usepackage{fancyhdr}
\usepackage[hidelinks]{hyperref}

% ---------------------------------------------------------------- colours (as in the brochure)
\definecolor{navy}{HTML}{1C3A5E}
\definecolor{pblue}{HTML}{3B78B5}
\definecolor{amber}{HTML}{F2B33D}
\definecolor{deepamber}{HTML}{D9902A}
\definecolor{pale}{HTML}{EEF3F9}
\definecolor{ink}{HTML}{2A2F36}
\definecolor{logonavy}{HTML}{173D63}
\definecolor{logodark}{HTML}{C98A1E}
\definecolor{logoamber}{HTML}{F2B446}
\definecolor{logocream}{HTML}{FEEFD2}
\definecolor{soft}{HTML}{5E6773}
\definecolor{agree}{HTML}{2F7D4F}
\definecolor{differ}{HTML}{B23B3B}
\color{ink}

\setlength{\parindent}{0pt}
\setlength{\parskip}{4pt}
\setlist{nosep,leftmargin=*,itemsep=2pt}
\renewcommand{\arraystretch}{1.15}

\pagestyle{fancy}
\fancyhf{}
\renewcommand{\headrulewidth}{0pt}
\fancyfoot[L]{\footnotesize\color{soft}\textbf{Pollis}\enspace New Methods\enspace·\enspace\nmtheme}
\fancyfoot[R]{\footnotesize\color{soft}\thepage}
\newcommand{\nmtheme}{}
\newcommand{\vignetteauthor}{}
\newcommand{\vignettedate}{\today}
\newcommand{\vignettelicense}{\href{https://creativecommons.org/licenses/by-sa/4.0/}{CC BY-SA 4.0}}

% The Pollis logo, redrawn so it prints sharply at any size
\newcommand{\pollislogo}[1]{%
  \begin{tikzpicture}[x=#1/512,y=-#1/512,baseline=(current bounding box.center)]
    \fill[logonavy] (0,0) rectangle (512,512);
    \fill[logodark] (36,422) -- (104,256) -- (172,422) -- cycle;
    \fill[logoamber] (104,422) -- (203,180) -- (302,422) -- cycle;
    \fill[logocream] (203,422) -- (339,89) -- (476,422) -- cycle;
  \end{tikzpicture}}

% Top of the first page: logo, name; then the kicker, title and a one-line promise
\newcommand{\nmhead}[3]{%
  \renewcommand{\nmtheme}{#1}%
  \thispagestyle{fancy}%
  \begin{minipage}[c]{13mm}\pollislogo{11mm}\end{minipage}%
  \begin{minipage}[c]{\dimexpr\linewidth-13mm\relax}
    {\color{navy}\Large\bfseries Pollis}\enspace{\color{soft}New Methods}\hfill
    {\color{deepamber}\small\bfseries\MakeUppercase{#1}}
  \end{minipage}\par\vspace{7mm}
  {\color{navy}\fontsize{22}{26}\selectfont\bfseries #2}\par\vspace{4pt}
  {\color{soft}\large #3}\par\vspace{5mm}}

\newcommand{\nmsection}[1]{\par\vspace{5pt}{\color{navy}\large\bfseries #1}\par\vspace{1pt}}
\newcommand{\nmsub}[1]{\par\vspace{2pt}{\color{deepamber}\bfseries #1}\par}
\newcommand{\gt}{\,{\color{deepamber}\bfseries ›}\,}
\newcommand{\menupath}[1]{{\bfseries\color{navy}#1}}
\newcommand{\fn}[1]{\texttt{\small #1}}
% Citations: \nmcite{key}{Breiman (2001)} links to the entry \nmref{key}{...} in the References
\newcommand{\nmcite}[2]{\hyperlink{ref:#1}{\color{pblue}#2}}
\newcommand{\nmref}[2]{\raisebox{1.1em}[0pt][0pt]{\hypertarget{ref:#1}{}}#2\par}
\newcommand{\nmdoi}[1]{\href{https://doi.org/#1}{\color{pblue}doi:#1}}
% Where the package is: its name and a link (its GitHub repository, once published)
\newcommand{\package}[2]{\href{#2}{\fn{#1.jl}}}

% Key facts at the top: source, missing in Julia, package, result
\newcommand{\factsheet}[4]{%
  \begin{tcolorbox}[enhanced, colback=pale, colframe=pale, boxrule=0pt, sharp corners,
    left=7pt, right=7pt, top=5pt, bottom=5pt, fontupper=\small]
  \begin{tabularx}{\linewidth}{@{}>{\leavevmode\bfseries\color{navy}}p{27mm}>{\leavevmode}X@{}}
  Tested against & #1\\
  Missing in Julia & #2\\
  Package & #3\\
  Result & #4
  \end{tabularx}
  \end{tcolorbox}}

% Boxes
\newtcolorbox{pbox}[2][]{enhanced, colback=pale, colframe=amber, boxrule=0pt,
  leftrule=3pt, sharp corners, left=7pt, right=7pt, top=6pt, bottom=6pt,
  fontupper=\small\raggedright, before upper={{\color{navy}\normalsize\bfseries #2}\par\smallskip}, #1}
\newtcolorbox{nbox}[2][]{enhanced, colback=navy, colframe=navy, boxrule=0pt,
  sharp corners, left=8pt, right=8pt, top=7pt, bottom=7pt, coltext=white,
  fontupper=\small\raggedright, before upper={{\color{amber}\normalsize\bfseries #2}\par\smallskip}, #1}

% Value tables: published beside Julia
\newenvironment{valuetable}[1]{%
  \par\small\tabularx{\linewidth}{#1}\toprule}{\bottomrule\endtabularx\par}
\newcommand{\ok}{{\color{agree}\faCheck}}
\newcommand{\no}{{\color{differ}\faTimes}}

% End of the document: credit and licence
\newcommand{\nmcredit}{%
  \par\vfill
  \begin{minipage}[b]{0.78\linewidth}\footnotesize\color{soft}
  A Pollis method, tested against the published tables on the published data.\par
  \textcopyright\ \vignetteauthor\enspace·\enspace Licensed under \vignettelicense\enspace·\enspace\vignettedate
  \end{minipage}\hfill
  \begin{minipage}[b]{0.18\linewidth}\raggedleft\pollislogo{9mm}\end{minipage}}


% A package Pollis needs for a panel of its own, not one that reproduces a paper: its own header, key facts and credit
\renewcommand{\vignetteauthor}{Antonio Saragga Seabra}
\renewcommand{\vignettedate}{2 November 2026}
\renewcommand{\vignettelicense}{MIT}
\fancyfoot[L]{\footnotesize\color{soft}\textbf{Pollis}\enspace Domain Packages\enspace·\enspace\nmtheme}
\renewcommand{\nmhead}[3]{%
  \renewcommand{\nmtheme}{#1}%
  \thispagestyle{fancy}%
  \begin{minipage}[c]{13mm}\pollislogo{11mm}\end{minipage}%
  \begin{minipage}[c]{\dimexpr\linewidth-13mm\relax}
    {\color{navy}\Large\bfseries Pollis}\enspace{\color{soft}Domain Packages}\hfill
    {\color{deepamber}\small\bfseries\MakeUppercase{#1}}
  \end{minipage}\par\vspace{7mm}
  {\color{navy}\fontsize{22}{26}\selectfont\bfseries #2}\par\vspace{4pt}
  {\color{soft}\large #3}\par\vspace{5mm}}
\renewcommand{\factsheet}[4]{%
  \begin{tcolorbox}[enhanced, colback=pale, colframe=pale, boxrule=0pt, sharp corners,
    left=7pt, right=7pt, top=5pt, bottom=5pt, fontupper=\small]
  \begin{tabularx}{\linewidth}{@{}>{\leavevmode\bfseries\color{navy}}p{27mm}>{\leavevmode}X@{}}
  In Pollis & #1\\
  Missing in Julia & #2\\
  Package & #3\\
  Tests & #4
  \end{tabularx}
  \end{tcolorbox}}
\renewcommand{\nmcredit}{%
  \par\vfill
  \begin{minipage}[b]{0.78\linewidth}\footnotesize\color{soft}
  Package written and tested with Claude Code (Anthropic) in Pollis.\par
  \textcopyright\ 2026 \vignetteauthor\enspace·\enspace Licensed under \vignettelicense\enspace·\enspace\vignettedate
  \end{minipage}\hfill
  \begin{minipage}[b]{0.18\linewidth}\raggedleft\pollislogo{9mm}\par\vspace{2pt}{\footnotesize contact@pollis.org}\end{minipage}}
\newcolumntype{R}{>{\raggedleft\arraybackslash}X}

\begin{document}
\emergencystretch=3em

\nmhead{Volatility}
  {StochasticVolatility: Heston, Bates and SABR in Julia}
  {Simulation, closed forms, the Hagan formula, Longstaff-Schwartz and finite differences (MCS-IT) for European
   and American prices}

\factsheet
  {Written to power the \emph{Continuous-Time Stochastic Volatility Models} panel of the Finance Toolbox. The panel
   does not call it yet.}
  {American prices under Bates and SABR. Of the packages read: DiffEqFinancial 2.11.0 has a Heston SDE and moments, no
   prices; Hedgehog 0.1.0 is Heston only, with no dividend yield (read from its source, not run); ArbitrageFreeSABR
   1.0.0 has the arbitrage-free SABR PDE, not the 2002 formula; QuantLib.jl 0.1.1 does not install on Julia 1.13, and
   its \fn{BatesEngine} at commit fe98714 has \fn{expm1(nu * delta2)} where the C++ has \fn{exp(nu + delta2) - 1}.}
  {\package{StochasticVolatility}{https://github.com/Trumpingtons/StochasticVolatility.jl} 0.1.0, MIT:\newline
   \fn{Pkg.Registry.add(url = "https://github.com/Trumpingtons/PollisRegistry.git")}\newline \fn{Pkg.add("StochasticVolatility")}}
  {\fn{test/runtests.jl} runs on Julia 1.10 and the latest 1.x at every push. Of 39 published American puts under
   Heston, 36 to every printed digit. Of 5 American calls under Bates, 3 agree and 2 are open. Of 14 American options
   under SABR, all within .00097 of QuantLib.}

\nmsection{Why a package}
The panel's procedures are hand-written loops, with no American price. American options have no closed form, so one
engine alone cannot show that it is right. The package has two: least squares on
simulated paths (\nmcite{ls2001}{Longstaff and Schwartz, 2001}) and finite differences on the pricing equation.

\nmsection{What it computes}
\fn{Heston}: $dS = \mu S\,dt + \sqrt{v}\,S\,dW_1$, $dv = \kappa(\theta - v)\,dt + \xi\sqrt{v}\,dW_2$, $d\langle W_1, W_2\rangle = \rho\,dt$
(\nmcite{heston1993}{Heston, 1993}). \fn{Bates} adds compensated log-normal jumps of intensity $\lambda$, with log size
of mean $\mu_J$ and deviation $\delta_J$ (\nmcite{bates1996}{Bates, 1996}). \fn{SABR}: $dF = \sigma F^\beta\,dW_1$,
$d\sigma = \nu\sigma\,dW_2$, $\sigma_0 = \alpha$ (\nmcite{hagan2002}{Hagan et al., 2002}). The drift $\mu$ is the rate less
the dividend yield.

\nmsub{Simulation}
\fn{simulate} and \fn{simulate\_terminal} with \fn{QuadraticExponential()} (\nmcite{andersen2008}{Andersen, 2008}),
\fn{FullTruncation()} (\nmcite{lord2010}{Lord et al., 2010}) or \fn{Euler()}, on \fn{PseudoRandom} or \fn{QuasiRandom} numbers:
scrambled Sobol with a Brownian bridge, and for Bates the jump bridge of \nmcite{baldeaux2012}{Baldeaux and Roberts (2012)}.

\nmsub{European prices}
\fn{european\_price}, the closed forms of Heston and Bates; \fn{hagan\_implied\_vol} and \fn{hagan\_price}, the formula of
Hagan et al., an approximation; \fn{monte\_carlo\_price}, a vector of strikes from one simulation.

\nmsub{American prices}
\fn{longstaff\_schwartz\_price} exercises at the simulation's steps, so it is Bermudan with as many dates.
By finite differences, \fn{finite\_difference\_price} is the MCS-IT scheme of \nmcite{hh2015}{Haentjens and in 't Hout (2015)}: modified
Craig-Sneyd ADI with $\theta = 1/3$ and the Ikonen-Toivanen step (\nmcite{it2009}{Ikonen and Toivanen, 2009}) for early
exercise, on nonuniform meshes with first-order upwind convection. For Bates the jump integral is explicit, exact on each
mesh interval under the log-normal jump. For SABR it works in $F$ and $\log\sigma$ with central differences, and
\fn{american = false} gives the European price. The default grid is 300 prices, 150 variances (or volatilities) and 75 steps.

\nmsection{A worked example}
\begin{footnotesize}
\begin{verbatim}
using StochasticVolatility
model = Heston(S0 = 10, v0 = 0.0625, mu = 0.1, kappa = 5, theta = 0.16, xi = 0.9, rho = 0.1)
finite_difference_price(model, 10, 0.25; rate = 0.1, call = false)   # 0.5204, American
european_price(model, 10, 0.25; rate = 0.1, call = false)            # 0.5015, European
sabr = SABR(100.0, 2.0, 0.5, -0.7, 0.8)                              # F0, alpha, beta, rho, nu
finite_difference_price(sabr, 100, 0.25; rate = 0.05, call = false)  # 3.9466, American
hagan_price(sabr, 100, 0.25; rate = 0.05, call = false)              # 3.9388, European
\end{verbatim}
\end{footnotesize}

\newpage
\nmsection{What the tests show: American puts under Heston, against the published MCS-IT}
From \nmcite{hh2015}{Haentjens and in 't Hout (2015)}, $K = 10$, $T = .25$ in Tables 2 and 3 and $K = 100$, $T = .25$ in Table 4. The
rows show the finest grid of each table; the tests compare all 39 values, on the same method and grids.

\begin{valuetable}{@{}l l *{5}{R} c@{}}
Source and grid & $S_0$ & 8 & 9 & 10 & 11 & 12 & \\ \midrule
Table 2, $v_0 = .0625$ & published & 2.0000 & 1.1081 & .5204 & .2143 & .0827 & \\
300 $\times$ 150 $\times$ 75 & Julia & 2.0000 & 1.1081 & .5204 & .2143 & .0827 & \ok \\ \midrule
Table 3, $v_0 = .25$ & published & 2.0788 & 1.3339 & .7962 & .4486 & .2433 & \\
300 $\times$ 150 $\times$ 75 & Julia & 2.0788 & 1.3339 & .7962 & .4486 & .2433 & \ok \\ \midrule
 & $S_0$ & 90 & 100 & 110 & & & \\ \midrule
Table 4, $v_0 = .0348$ & published & 10.0039 & 3.2126 & .9305 & & & \\
300 $\times$ 150 $\times$ 60 & Julia & 10.0039 & 3.2126 & .9305 & & & \ok \\
\end{valuetable}

\medskip
\noindent\textbf{Where they differ.} 36 of 39 agree to every printed digit, the 3 others at the exercise boundary, with
few steps or points. $S_0 = 8$ at 100 $\times$ 50 $\times$ 25: published 2.0001, Julia 2.0000, off by 1 in the fourth decimal.
$S_0 = 90$ at 20 steps: published 9.9984, Julia 10.0000, off by 16; the published value is below the exercise value 10.
$S_0 = 90$ at 40 steps: published 10.0015, Julia 10.0020, off by 5. Of the 33 values not within a cent of the exercise
value, the furthest is $4.9\times10^{-5}$ away, under half of the last digit printed.

\nmsub{American calls under Bates: 2 of 5 open}
$K = 100$, $T = .5$, rate .03, dividend yield .05, $v_0 = \theta = .04$, $\kappa = 2$, $\xi = .25$, $\rho = -.5$, $\lambda = .2$, $\mu_J = -.58$,
$\delta_J = .4$. Published: \nmcite{toivanen2010}{Toivanen (2010)} on a grid of 4096 $\times$ 2048 $\times$ 514, read second-hand from
\nmcite{salmi2013}{Salmi, Toivanen and von Sydow (2013)}. The last column is the value extrapolated from the two grids of
Julia, $2V_{1200} - V_{600}$, minus the published one. The 1200 grid, QuantLib 1.44 and Longstaff-Schwartz are stored in
\fn{benchmark/}, not recomputed here.

\begin{valuetable}{@{}c *{5}{R} r c@{}}
$S_0$ & Published & Julia 600 $\times$ 300 $\times$ 150 & Julia 1200 $\times$ 600 $\times$ 300 & QuantLib 100 $\times$ 100 $\times$ 50 & LS, 98304 paths, 126 steps (SE) & Extrap. $-$ publ. & \\ \midrule
80 & .328526 & .332788 & .330687 & .331071 & .33142 (.00174) & +.0001 & \ok \\
90 & 2.109397 & 2.115783 & 2.112723 & 2.108276 & 2.09304 (.00330) & +.0003 & \ok \\
100 & 6.711622 & 6.715916 & 6.713882 & 6.714004 & 6.69703 (.00723) & +.0002 & \ok \\
110 & 13.749337 & 13.749886 & 13.748658 & 13.759956 & 13.70704 (.00850) & $-$.0019 & \no \\
120 & 22.143307 & 22.138758 & 22.137875 & 22.153021 & 22.08472 (.01164) & $-$.0063 & \no \\
\end{valuetable}

\medskip
\nmsub{American options under SABR, against QuantLib}
$F_0 = 100$, $\alpha = 2$, $\beta = .5$, $\rho = -.7$, $\nu = .8$, $T = .25$, rate .05. QuantLib 1.44 \fn{FdSabrVanillaEngine} and Julia at
400 $\times$ 50 $\times$ 64, the largest gap .00097 at call $K = 75$. LS (stored, with SE): 131072 quasi-random paths, 64 steps.

\begin{valuetable}{@{}c *{6}{R}@{}}
$K$ & Put, QuantLib & Put, Julia & Put, LS & Call, QuantLib & Call, Julia & Call, LS \\ \midrule
50 & .00244 & .00251 & .00271 (.00029) & 50 & 50 & 50 \\
75 & .12580 & .12598 & .12513 (.00068) & 25.02978 & 25.03075 & 25.03815 (.00490) \\
95 & 2.15536 & 2.15524 & 2.14291 (.00275) & 7.11113 & 7.11123 & 7.10076 (.00475) \\
100 & 3.94654 & 3.94666 & 3.93968 (.00368) & 3.94646 & 3.94650 & 3.94508 (.00264) \\
105 & 6.72499 & 6.72566 & 6.71495 (.00508) & 1.76876 & 1.76890 & 1.77113 (.00222) \\
125 & 25 & 25 & 25 & .00276 & .00217 & .00338 (.00015) \\
200 & 100 & 100 & 100 & 0 & 0 & 0 \\
\end{valuetable}

\medskip
\nmsub{Cases with a known answer}
An American call with no dividend is never exercised early, so it must tend to the closed form. The gap below halves as
the grid doubles: first order, largest for $S_0$ below $K$. $S_0 = 100$, $v_0 = \theta = .04$, rate .05, $T = .5$, $\kappa = 2$, $\xi = .3$, $\rho = -.7$;
Bates with $\lambda = .2$, $\mu_J = -.58$, $\delta_J = .4$ (\fn{benchmark/finite\_difference\_no\_dividend\_call.txt}).

\begin{valuetable}{@{}l c *{4}{R}@{}}
Model & $K$ & Closed form & Gap at 300 $\times$ 150 $\times$ 75 & at 600 $\times$ 300 $\times$ 150 & at 1200 $\times$ 600 $\times$ 300 \\ \midrule
Heston & 80 & 22.436161 & .001693 & .000861 & .000434 \\
Heston & 100 & 6.837200 & .006962 & .003546 & .001789 \\
Heston & 120 & .504430 & .013094 & .006700 & .003388 \\
Bates & 80 & 24.325572 & .003762 & .001886 & \\
Bates & 100 & 8.960508 & .016100 & .008099 & \\
Bates & 120 & 1.147607 & .042395 & .021430 & \\
\end{valuetable}

\medskip
\noindent Under SABR at the 400 $\times$ 50 $\times$ 64 grid, American minus European at rate 0 is $3.8\times10^{-9}$, $1.2\times10^{-6}$ and
$6.7\times10^{-5}$ at $K = 75$, 100 and 125, and put-call parity holds to about $10^{-13}$. At $\nu = .001$, $K = 100$, Julia gives 3.939869
with $\beta = 0$ against Bachelier's 3.939865, and 3.938194 with $\beta = 1$ against Black's 3.938224. Longstaff-Schwartz under
Heston (131072 quasi-random paths) is within .005 of the published values at 64 dates, and within .01 of the published
Bermudan values at 60.

\nmsub{European prices and simulation}
\begin{valuetable}{@{}l r l l c@{}}
Function & Rows & Against & Largest difference & \\ \midrule
\fn{european\_price}, Heston & 189 & QuantLib 1.44 / NMOF 2.12.0 & $2.8\times10^{-12}$ / $9.0\times10^{-7}$ & \ok \\
\fn{european\_price}, Bates & 540 & QuantLib, 324 rows with $\xi < 1$ / NMOF & $1.4\times10^{-8}$ / $1.1\times10^{-6}$ & \ok \\
\fn{hagan\_implied\_vol}, \fn{hagan\_price} & 756 & QuantLib 1.44 & $1.5\times10^{-14}$, relative & \ok \\
\end{valuetable}

\medskip
\noindent QuantLib's \fn{BatesEngine} has not converged at $\xi \ge 1$, so there only NMOF is used. On one file of Sobol numbers,
16 replicates of $2^{13}$ paths and 64 steps, the paths agree with QuantLib to $9.1\times10^{-12}$ or closer, at 4.5 against 27.0 µs per path for
Heston and 6.0 to 6.2 against 30.2 for Bates. QuantLib has no SABR process, so its side is written by hand on its generator: 4.5 to 4.6
against 4.0 µs per path, with Euler. \fn{QuasiRandom} narrows the spread of replicate means 12 to 21 times for Heston and SABR calls, and 2.8
to 5.0 times for Bates calls.

\medskip
\noindent One call price, in µs: Julia 155.1 (Heston) and 176.8 (Bates), \fn{rtol} $10^{-12}$; QuantLib C++ with Laguerre order 192, 1907.0 and
1947.4, and adaptive $10^{-13}$, 678.8 for Heston; NMOF 1460.3 and 1763.0. \fn{hagan\_price}: .164 against .104. The 14 SABR options
on one thread: .66 to .73 s against 3.05 s. These are not times at equal accuracy.

\nmsection{Still open}
\begin{itemize}[leftmargin=*, itemsep=1pt, topsep=2pt]
\item Calibration is not written.
\item Greeks are not written.
\end{itemize}

\nmsection{References}
{\footnotesize\raggedright\setlength{\parskip}{1pt}
\nmref{hh2015}{Haentjens, T. and in 't Hout, K. J. (2015). ADI schemes for pricing American options under the Heston model. \emph{Applied Mathematical Finance}, 22(3), 207--237. \nmdoi{10.1080/1350486X.2015.1009129}. Values read from arXiv:1309.0110v1.}
\nmref{andersen2008}{Andersen, L. (2008). Simple and efficient simulation of the Heston stochastic volatility model. \emph{Journal of Computational Finance}, 11(3), 1--42. \nmdoi{10.21314/JCF.2008.189}}
\nmref{baldeaux2012}{Baldeaux, J. and Roberts, D. (2012). Quasi-Monte Carlo methods for the Heston model. arXiv:1202.3217.}
\nmref{bates1996}{Bates, D. S. (1996). Jumps and stochastic volatility: exchange rate processes implicit in Deutsche Mark options. \emph{Review of Financial Studies}, 9(1), 69--107. \nmdoi{10.1093/rfs/9.1.69}}
\nmref{hagan2002}{Hagan, P. S., Kumar, D., Lesniewski, A. S. and Woodward, D. E. (2002). Managing smile risk. \emph{Wilmott Magazine}, 84--108.}
\nmref{heston1993}{Heston, S. L. (1993). A closed-form solution for options with stochastic volatility with applications to bond and currency options. \emph{Review of Financial Studies}, 6(2), 327--343. \nmdoi{10.1093/rfs/6.2.327}}
\nmref{it2009}{Ikonen, S. and Toivanen, J. (2009). Operator splitting methods for pricing American options under stochastic volatility. \emph{Numerische Mathematik}, 113(2), 299--324. \nmdoi{10.1007/s00211-009-0227-5}}
\nmref{ls2001}{Longstaff, F. A. and Schwartz, E. S. (2001). Valuing American options by simulation: a simple least-squares approach. \emph{Review of Financial Studies}, 14(1), 113--147. \nmdoi{10.1093/rfs/14.1.113}}
\nmref{lord2010}{Lord, R., Koekkoek, R. and van Dijk, D. (2010). A comparison of biased simulation schemes for stochastic volatility models. \emph{Quantitative Finance}, 10(2), 177--194. \nmdoi{10.1080/14697680802392496}}
\nmref{salmi2013}{Salmi, S., Toivanen, J. and von Sydow, L. (2013). Iterative methods for pricing American options under the Bates model. \emph{Procedia Computer Science}, 18, 1136--1144. \nmdoi{10.1016/j.procs.2013.05.279}}
\nmref{toivanen2010}{Toivanen, J. (2010). A componentwise splitting method for pricing American options under the Bates model. In \emph{Applied and Numerical Partial Differential Equations}, Computational Methods in Applied Sciences, Springer, 213--227. \nmdoi{10.1007/978-90-481-3239-3\_16}}
}

\nmcredit
\end{document}
